% !PRINTED 108
\weberpage{108}
%<<CONT

$\xi$. The same conclusion may be drawn in quite the same way when $f(a)$ is positive and $f(b)$ negative, and so our theorem is proved.

From this we draw two important consequences. If $f(\xi)$ vanishes, then an interval
\[
 \xi-\varepsilon<x<\xi+\varepsilon
\]
can be so determined that $f(x)$ does not vanish in this interval except for $x=\xi$. Then four cases are possible with respect to the signs of $f(x)$ in the two intervals $\xi-\varepsilon<x<\xi$ and $\xi<x<\xi+\varepsilon$, which are put together in the following scheme:
\[
\begin{array}{lcc}
&\multicolumn{2}{c}{\xi-\varepsilon<\xi<\xi+\varepsilon}\\
1.\ f(x)&+&-\\
2.\ f(x)&-&+\\
3.\ f(x)&+&+\\
4.\ f(x)&-&-.
\end{array}
\]
If $f^{(\nu)}(x)$ is the first among the derivatives of $f(x)$ which is different from zero for $x=\xi$, then the formula [§ 13, (2)]
\[
 f(\xi+h)=\frac{h^\nu}{\Pi(\nu)}f^{(\nu)}(\xi)+\cdots,
\]
shows that these four cases are distinguished by the following signs:
\[
\begin{array}{lll}
1.&\nu\text{ odd}& f^{(\nu)}(\xi)<0\\
2.&\nu\text{ odd}& f^{(\nu)}(\xi)>0\\
3.&\nu\text{ even}& f^{(\nu)}(\xi)>0\\
4.&\nu\text{ even}& f^{(\nu)}(\xi)<0
\end{array}
\]
We say that in the first case $f(x)$ passes \gesperrt{decreasing} through zero, and in the second \gesperrt{increasing} through zero. In case 3, $f(\xi)=0$ is a \gesperrt{minimum}, in case 4 a \gesperrt{maximum}.

In particular, if $f'(\xi)$ is different from zero, so that $\xi$ is a simple root, then the positive or negative sign of $f'(\xi)$ is the characteristic mark for the increase or decrease of the function $f(x)$ as $x$ passes through $\xi$.

With these aids we still prove two further important theorems.

\medskip
\noindent IX. \gesperrt{The equation}\par
\[
 x^n-a=0
\]
\gesperrt{has, when $a$ is a positive number, one and only one positive root.}\par
\medskip
